\documentclass[12pt]{article}
\usepackage[utf8]{inputenc}
\usepackage[T1]{fontenc}
\usepackage{hyperref}
\title{Talondo' Language Resources for K-12 Curriculum Development}
\author{Compiled by Claude Sonnet 5 (claude-sonnet-5) for Andr\'as Kornai}
\date{}
\begin{document}
\maketitle

\section{Introduction}
Talondo' (ISO 639-3 tln) is an endangered Austronesian language of the South Sulawesi branch (Northern Toraja group). It is spoken in Talondo Kondo village, Mamuju Regency, West Sulawesi, Indonesia. Speaker counts in the sources I found disagree: about 400 (SIL 2004, via endangeredlanguages.com), about 500 (Global Recordings Network) and 1,200 (Wikipedia, 2011). The language is under pressure from larger neighbouring languages such as Mamuju, and from Indonesian. The searches found no school materials in Talondo'.

\section{Primary Educational Resources}
The searches found no primers, textbooks or other school materials in Talondo'.

\section{Linguistic and Reference Materials}
\begin{itemize}
\item Wikipedia overview of the language: \url{https://en.wikipedia.org/wiki/Talondo%CA%BC_language}
\item Verbix endangered-languages entry: \url{https://docs.verbix.com/EndangeredLanguages/Talondo}
\item Endangered Languages Project entry: \url{https://endangeredlanguages.com/lang/2788}
\end{itemize}

\section{Religious and Community Resources}
No scripture or community organisation was found in the searches.

\section{Digital and Audio Resources}
\begin{itemize}
\item Global Recordings Network language page (I did not check what recordings it lists): \url{https://globalrecordings.net/en/language/tln}
\end{itemize}

\section{Curriculum Development Context}
Very little documentation was found online. A curriculum would need new fieldwork with the Talondo' community, working alongside Indonesian-medium schooling. The reference sources above are overviews only and are not teaching materials.
\end{document}